\chapterbackmatter{Supplementary Exercises}

\begin{enumerate}
\SupplementaryExercise{1}{A Census of One-Loop QED Graphs}
  {(Adapted from Bernhard Mistlberger and Kevin Zhou, Stanford PHYS 330,
  Problem Set 9, Problem 1)}
  {stanford-phys330-ps9-2022-ch11}
For \(e^-e^+\to\mu^-\mu^+\), classify every connected diagram of order
\(e^4\).  Separate external-leg, vacuum-polarization, vertex, and box
topologies, and explain which pieces are removed or modified by on-shell
renormalization.

\SupplementaryExercise{2}{The General Feynman-Parameter Simplex}
  {(Adapted from Bernhard Mistlberger and Kevin Zhou, Stanford PHYS 330,
  Problem Set 9, Problem 2)}
  {stanford-phys330-ps9-2022-ch11}
Starting from the Schwinger representation for \(A_i^{-\nu_i}\), derive the
general \(n\)-denominator Feynman-parameter identity with arbitrary positive
\(\nu_i\).  Include the simplex delta function and all Gamma-function
normalizations.

\SupplementaryExercise{3}{The Massless Bubble in \(d\) Dimensions}
  {(Adapted from Bernhard Mistlberger and Kevin Zhou, Stanford PHYS 330,
  Problem Set 9, Problem 2)}
  {stanford-phys330-ps9-2022-ch11}
Evaluate
\[
 I_B(a,b;p^2)=\int\frac{d^dk}{(2\pi)^d}
 \frac{1}{(k^2-i0)^a((k+p)^2-i0)^b}
\]
by Feynman parametrization and Wick rotation.  Give the answer for general
\(a,b,d\), and identify the ultraviolet pole of \(I_B(1,1;p^2)\) near
\(d=4\).

\SupplementaryExercise{4}{Rank-Four Rotational Averaging}
  {(Adapted from Bernhard Mistlberger and Kevin Zhou, Stanford PHYS 330,
  Problem Set 9, Problem 3)}
  {stanford-phys330-ps9-2022-ch11}
For a rotationally invariant function \(f(k^2)\) in \(d\) Euclidean
dimensions, reduce the integrals with numerators \(k^\mu k^\nu k^\rho\)
and \(k^\mu k^\nu k^\rho k^\sigma\) to scalar integrals.  Determine every
coefficient by contraction.

\SupplementaryExercise{5}{Tensor Reduction of a Bubble}
  {(Adapted from Bernhard Mistlberger and Kevin Zhou, Stanford PHYS 330,
  Problem Set 9, Problem 3)}
  {stanford-phys330-ps9-2022-ch11}
Reduce the massless rank-one and rank-two bubble integrals to the scalar
bubble \(I_B(a,b;p^2)\) and tadpoles.  In dimensional regularization, use the
vanishing of scaleless tadpoles to obtain the particularly simple result for
\(a=b=1\).

\SupplementaryExercise{6}{Massless QED Self-Energy Poles}
  {(Adapted from Bernhard Mistlberger and Kevin Zhou, Stanford PHYS 330,
  Problem Set 9, Problem 4)}
  {stanford-phys330-ps9-2022-ch11}
In Feynman gauge and \(d=4-2\varepsilon\), extract the ultraviolet poles of
the massless one-loop electron self-energy and photon vacuum polarization.
Show explicitly that the photon pole is transverse.

\SupplementaryExercise{7}{The Area of a \(d\)-Sphere}
  {(Adapted from Heidelberg University, Quantum Field Theory 1, Sheet 11,
  Problem 1(a))}
  {heidelberg-qft1-sheet11-2009-ch11}
Evaluate the \(d\)-dimensional Gaussian both as a product of one-dimensional
integrals and in radial coordinates.  Derive
\(\Omega_d=2\pi^{d/2}/\Gamma(d/2)\) and state its analytic continuation in
\(d\).

\SupplementaryExercise{8}{A Massive Euclidean Master Integral}
  {(Adapted from Heidelberg University, Quantum Field Theory 1, Sheet 11,
  Problem 1(b))}
  {heidelberg-qft1-sheet11-2009-ch11}
Evaluate
\[
 J_n(m^2)=\int\frac{d^dq}{(2\pi)^d}\frac{1}{(q^2+m^2)^n}
\]
for complex \(d\) and \(n\).  State its convergence domain before analytic
continuation and locate its ultraviolet poles.

\SupplementaryExercise{9}{Completing a Shifted Square}
  {(Adapted from Heidelberg University, Quantum Field Theory 1, Sheet 11,
  Problem 1(c))}
  {heidelberg-qft1-sheet11-2009-ch11}
Use a justified momentum shift to evaluate
\[
 \int\frac{d^dq}{(2\pi)^d}
 \frac{1}{(q^2+m^2+2q\cdot p)^n}.
\]
Specify the condition on \(m^2-p^2\) in the Euclidean domain and then give
the analytically continued result.

\SupplementaryExercise{10}{Shifted Vector and Tensor Integrals}
  {(Adapted from Heidelberg University, Quantum Field Theory 1, Sheet 11,
  Problem 1(c))}
  {heidelberg-qft1-sheet11-2009-ch11}
Evaluate the shifted master integrals of Exercise S.11.9 with numerators
\(q^\mu\) and \(q^\mu q^\nu\).  Express the answers in terms of \(p^\mu\),
\(\delta^{\mu\nu}\), and Gamma functions, and check them by differentiating
the scalar master integral.

\SupplementaryExercise{11}{The Spinor Vacuum-Polarization Numerator}
  {(Adapted from Heidelberg University, Quantum Field Theory 1, Sheet 11,
  Problem 2(a)--(b))}
  {heidelberg-qft1-sheet11-2009-ch11}
Starting from the closed-fermion-loop rules, take the Dirac trace in the
one-loop photon self-energy, introduce one Feynman parameter, and shift the
loop momentum.  Show before integration how the terms organize into a
transverse tensor plus a rotational average.

\SupplementaryExercise{12}{Subtracted Vacuum Polarization at Small Momentum}
  {(Adapted from Heidelberg University, Quantum Field Theory 1, Sheet 11,
  Problem 2(c))}
  {heidelberg-qft1-sheet11-2009-ch11}
Renormalize the spinor-QED vacuum polarization by imposing \(\pi(0)=0\).
Derive the finite parameter integral for \(\pi(q^2)\) and its first two
terms for \(\lvert q^2\rvert\ll m^2\).

\SupplementaryExercise{13}{Which QED Counterterms Are Allowed?}
  {(Adapted from Mark Srednicki, Quantum Field Theory, Chapter 62,
  Eqs. (62.1)--(62.6))}
  {srednicki-qft-qed-loops-ch11}
Use Lorentz invariance, gauge invariance, charge conjugation, parity, and
power counting to list all local QED counterterms of mass dimension at most
four.  Explain why no photon mass, photon tadpole, or Pauli term belongs to
the renormalizable counterterm Lagrangian.

\SupplementaryExercise{14}{Photon Pole and Charge Normalization}
  {(Adapted from Mark Srednicki, Quantum Field Theory, Chapter 62,
  Eqs. (62.7)--(62.25))}
  {srednicki-qft-qed-loops-ch11}
Dyson-sum a transverse photon self-energy and show how the residue at
\(q^2=0\) fixes the on-shell condition \(\pi(0)=0\).  Relate this condition
to the definition of the measured electric charge.

\SupplementaryExercise{15}{Two On-Shell Conditions for a Fermion}
  {(Adapted from Mark Srednicki, Quantum Field Theory, Chapter 62,
  Eqs. (62.26)--(62.39))}
  {srednicki-qft-qed-loops-ch11}
Write the most general parity-even fermion self-energy as
\(\Sigma(p)=A(p^2)i\sl p+B(p^2)m\).  Derive the two on-shell conditions that
fix the mass and field-strength counterterms, and show that the renormalized
self-energy has a double zero in the inverse propagator expansion about the
physical pole.

\SupplementaryExercise{16}{The Equality of Vertex and Wave-Function Poles}
  {(Adapted from Mark Srednicki, Quantum Field Theory, Chapters 62 and 68,
  especially Problem 62.2)}
  {srednicki-qft-qed-loops-ch11}
Use the zero-momentum Ward identity to prove that the ultraviolet pole of
the QED vertex counterterm equals that of the electron wave-function
counterterm.  State the result as \(Z_1=Z_2\) and explain what it implies for
charge renormalization.

\SupplementaryExercise{17}{Why Four-Photon Scattering Is Ultraviolet Finite}
  {(Adapted from Mark Srednicki, Quantum Field Theory, Chapter 62,
  Problem 62.3)}
  {srednicki-qft-qed-loops-ch11}
Each one-loop box ordering with four external photons is superficially
logarithmically divergent.  Use gauge invariance and the absence of an
allowed dimension-four four-photon counterterm to show that the complete
sum is ultraviolet finite.  Also explain why the three-photon amplitude
vanishes.

\SupplementaryExercise{18}{From the Pauli Form Factor to \(g-2\)}
  {(Adapted from Mark Srednicki, Quantum Field Theory, Chapters 63--64)}
  {srednicki-qft-qed-loops-ch11}
Decompose the on-shell QED vertex into Dirac and Pauli form factors.  Couple
it to a slowly varying magnetic field and prove that
\((g_e-2)/2=F_2(0)\).  Identify which normalization fixes \(F_1(0)\).

\SupplementaryExercise{19}{The One-Loop QED Beta Function}
  {(Adapted from David Skinner, Advanced Quantum Field Theory,
  Section 6.3.3)}
  {skinner-aqft-qed-renormalization-2018-ch11}
Starting from the pole part of the photon field counterterm in minimal
subtraction, derive the one-loop beta function for one Dirac fermion,
\(\beta(e)=e^3/(12\pi^2)\).  Keep the \(d\)-dimensional classical scaling
term until the pole cancellation is explicit.

\SupplementaryExercise{20}{Vacuum Screening and the Running Charge}
  {(Adapted from David Skinner, Advanced Quantum Field Theory,
  Section 6.3.4)}
  {skinner-aqft-qed-renormalization-2018-ch11}
Integrate the one-loop QED beta function between two scales.  Show that the
effective charge grows at short distance and explain the sign using the
polarization cloud around a test charge.

\SupplementaryExercise{21}{Decoupling below a Heavy Threshold}
  {(Adapted from David Skinner, Advanced Quantum Field Theory,
  Sections 6.3.3--6.3.4)}
  {skinner-aqft-qed-renormalization-2018-ch11}
Use the on-shell-subtracted vacuum-polarization integral to show that a
charged fermion of mass \(M\) contributes only
\(O(q^2/M^2)\) to processes with \(\lvert q^2\rvert\ll M^2\).  Reconcile
this decoupling with the mass-independent minimal-subtraction beta function.

\SupplementaryExercise{22}{The Euler--Heisenberg Operator Basis}
  {(Adapted from David Skinner, Advanced Quantum Field Theory,
  Section 7.2)}
  {skinner-aqft-qed-renormalization-2018-ch11}
At energies well below the electron mass, classify the independent
Lorentz- and parity-even operators built from four electromagnetic field
strengths and no derivatives.  Show that only two invariants are required
and give them both covariantly and in terms of \(\mathbf E,\mathbf B\).

\SupplementaryExercise{23}{Expanding the Proper-Time Determinant}
  {(Adapted from David Skinner, Advanced Quantum Field Theory,
  Sections 7.1--7.2)}
  {skinner-aqft-qed-renormalization-2018-ch11}
For a constant electromagnetic background, the renormalized one-loop
spinor determinant may be written in Schwinger proper time.  Expand it
through fourth order in the fields and recover the Euler--Heisenberg
interaction, including its numerical coefficients.

\SupplementaryExercise{24}{Low-Energy Light-by-Light Scaling}
  {(Adapted from David Skinner, Advanced Quantum Field Theory,
  Section 7.2)}
  {skinner-aqft-qed-renormalization-2018-ch11}
Use the Euler--Heisenberg interaction to determine the energy and
fine-structure-constant dependence of the photon--photon scattering
amplitude and total cross section for \(\omega\ll m_e\).  Explain why no
lower power of \(\omega/m_e\) is possible.

\SupplementaryExercise{25}{The Pauli Operator in a Background Field}
  {(Adapted from Hong Liu and MIT OpenCourseWare, 8.324 Lecture 15,
  Eqs. (1)--(4))}
  {mit-8324-f10-lecture15-ch11}
Show that the zero-momentum Pauli form factor is reproduced by a local
operator \(\bar\psi\sigma^{\mu\nu}F_{\mu\nu}\psi\).  Match its coefficient
to the on-shell vertex and verify its mass dimension and discrete
symmetries.

\SupplementaryExercise{26}{Nonrelativistic Reduction with a Pauli Term}
  {(Adapted from Hong Liu and MIT OpenCourseWare, 8.324 Lecture 15,
  Eqs. (5)--(11))}
  {mit-8324-f10-lecture15-ch11}
Split the Dirac spinor into large and small two-component fields in a slowly
varying electromagnetic background.  Eliminate the small field through
order \(1/m\), and derive the magnetic Hamiltonian including \(F_2(0)\).

\SupplementaryExercise{27}{Three Denominators on a Simplex}
  {(Adapted from Hong Liu and MIT OpenCourseWare, 8.324 Lecture 15,
  Eqs. (12)--(17))}
  {mit-8324-f10-lecture15-ch11}
Apply the three-denominator Feynman identity to the one-loop QED vertex.
For on-shell external electrons, perform the loop-momentum shift and derive
the common denominator in terms of three simplex parameters and \(q^2\).

\SupplementaryExercise{28}{Isolating the Pauli Numerator}
  {(Adapted from Hong Liu and MIT OpenCourseWare, 8.324 Lecture 15,
  Eqs. (18)--(22))}
  {mit-8324-f10-lecture15-ch11}
Use symmetric integration, gamma-matrix identities, and the external Dirac
equations to isolate the part of the one-loop vertex numerator proportional
to \(p^\mu+p'{}^\mu\).  Explain why this is the part needed for \(F_2\).

\SupplementaryExercise{29}{A Finite Euclidean Vertex Integral}
  {(Adapted from Hong Liu and MIT OpenCourseWare, 8.324 Lecture 15,
  Eqs. (22)--(24))}
  {mit-8324-f10-lecture15-ch11}
Wick rotate and evaluate the convergent loop integral
\[
 \int\frac{d^4k}{(2\pi)^4}\frac{1}{(k^2+\Delta)^3}.
\]
Use it to reduce the one-loop Pauli form factor to a finite simplex
integral, retaining a nonzero momentum transfer \(q\).

\SupplementaryExercise{30}{Schwinger's Parameter Integral}
  {(Adapted from Hong Liu and MIT OpenCourseWare, 8.324 Lecture 15,
  Eqs. (24)--(25))}
  {mit-8324-f10-lecture15-ch11}
Take \(q^2\to0\) in the simplex representation of the Pauli form factor and
perform every remaining parameter integral.  Derive
\(F_2(0)=\alpha/(2\pi)\) and \(g_e=2+\alpha/\pi\).
\end{enumerate}
